\documentclass{article}
\usepackage{amsmath,amstext,amssymb}
\usepackage[margin=1in]{geometry}

\begin{document}
{\large\bf Percentages and Rates}\hfill Name: \underline{\hspace{2in}}\\

Please work the following questions with those around you. You will be
submitting the activity on Friday for credit at the end of the class (on Friday), 
instead of a quiz this week. \\

\begin{enumerate}
\item Write each of the following as a percent:
    \begin{enumerate}
    \item $\frac{1}{4}$
    \item 0.02
    \item 2.35
    \end{enumerate}

\item The sales tax in a town is 6.5\%. How much will you pay on a \$140
purchase? How much will you pay if you have a 20\% off coupon? Does is matter if
the tax is applied before the coupon, or after? Why or why not? 

\vfill

\item Suppose a certain stock costs \$11.43 when the market opens, and \$12.14
when the market closes. What percent increase is this? 

\vfill\newpage

\item A basketball player scores on 40\% of 2-point field goal attempts, and on
30\% of 3-point field goal attempts. Find the player's overall field goal
percentage. 

\begin{enumerate}
\item What is the average of these two percentages? \underline{\hspace{1in}}
\item Suppose that the player attempted 200 2-point field goals and 100 3-point
field goals. How many 2-point shots did they make? How many 3-point shots? What
is their overall field goal average ($\frac{\text{total shots made}}{\text{total
shots attempted}}$)
\item Do we have enough information to really answer the question as it's
stated? Why or why not? 
\end{enumerate}

\vfill

\item Your car can drive 300 miles on 15 gallons of gasoline. How many miles
per gallon does your car get? 

\vspace{1in}

\begin{enumerate}
\item If you have the same car and it has 12 gallons of gasoline in it's tank,
how far can you drive?
\end{enumerate}

\vfill

\item A map scale indicates that $\frac{1}{2}$ inch on the map corresponds to 3
miles. How many miles apart are two cities which are $2\frac{1}{4}$ inches apart
on the map? 

\vfill 
\end{enumerate}
\end{document}
